% !TeX root = compte-rendu.tex
% !TeX spellcheck = fr


\section{Coalitions réalistes}
Dans cette partie, nous munissons les votants d'une relation d'ordre totale "est plus à gauche que". Nous considérons donc $(\V, \preccurlyeq)$.
Nous avons créé la notion de coalitions réalistes pour étudier des votants placés sur le spectre politique droite/gauche (à une seule dimension).

\begin{definition}
	Une coalition $C \in \mathcal{P}(\V)$ est dite \textit{réaliste} si
	\begin{equation}
		\forall (i, j) \in \V^2, \forall k \in \V,
		i \preccurlyeq k \preccurlyeq j \implies k \in C
	\end{equation}
	Cette notion serait une sorte de convexité non continue. \\
	Notons $\R$ l'ensemble des coalitions réalistes et $\R_i$ l'ensemble de celles contenant l'électeur $i$.
\end{definition}

\begin{exemple}
	Considérons $\V = \{G, C, D\}$ avec $G \preccurlyeq C \preccurlyeq D$ pour gauche, centre et droite. Dessinons le diagramme de Hasse de $(\R, \subseteq)$ :
	\begin{equation*}
		\begin{tikzcd}[column sep=0, row sep=1.5cm]
			&& \{G, C, D\} \drar[dash] \dlar[dash] && \\
			& \{G, C\} \drar[dash] \dlar[dash] && \{C, D\} \drar[dash] \dlar[dash] & \\
			\{G\} \arrow[dash]{drr} && \{C\} \dar[dash] && \{D\} \arrow[dash]{dll} \\
			&& \emptyset && \\
		\end{tikzcd}
	\end{equation*}
\end{exemple}

Pour $\V$ quelconque, si on lit le digramme de Hasse de haut en bas, on obtient $|\R| = \left( \sum_{k=1}^{|\V|} i \right) + 1$.

\begin{equation}
	 |\R| = \frac{N \left(N+1\right)}{2} + 1
\end{equation}

Considérons que $\V = [\![1;N]\!]$.
Soit $i \in \V$.
On a $\R_i = \left\{ [\![j;j']\!] \;|\; (j, j') \in \V^2, j \leqslant i \leqslant j' \right\} = \left\{ [\![j;j']\!] \;|\; (j, j') \in [\![1;i]\!] \times [\![i;N]\!] \right\}$.
Donc
\begin{equation}
	\forall i \in V, \
	|\R_i| = i (N-i+1)
\end{equation}

Nous pouvons ainsi redéfinir la décidabilité et les indices de Banzhaf.

\begin{definition}
	Posons $d_i^r$ la décidabilité réaliste d'un électeur $i$ :
	\begin{equation}
		d_i^r
		= \left| \left\{ C \in \G_i \cap \R \;|\; C \setminus \{i\} \notin \G \right\} \right|
	\end{equation}
	De même, posons les indices de Banzhaf réalistes normalisé et absolu :
	\begin{equation*}
		B_i^r = \frac{d_i^r}{\displaystyle \sum_{j=1}^{n} d_j^r}
		\qquad \qquad
		B_i'^r = \frac{d_i^r}{\displaystyle |\R_i|} = \frac{d_i^r}{i(N-i+1)}
	\end{equation*}
\end{definition}

La \hyperlink{section.5}{partie 5} donne un exemple où le calcul formel de $B_i'^r$ est possible.
Décrivons ici un algorithme qui permet de calculer ces deux indices.
On sait que $\R = \left\{\emptyset\right\} \cup \left\{ [\![j;j']\!] \;|\; (j, j') \in \V^2, j \leqslant j' \right\}$. \\

\begin{algorithm}[H]
	\caption{Calcul de l'ensemble des coalitions réalistes de $(E, \preccurlyeq)$}
	\Entree{Une liste L des éléments de l'ensemble E triés selon $\preccurlyeq$}
	\Sortie{Ensemble des coalitions réalistes}
	
	$N \gets Longueur(L)$ \\
	$R \gets \{ \emptyset \}$ \\
	
	\Pour{$i$ de $1$ à $N$}{
		\Pour{$j$ de $i$ à $N$}{
			Ajouter à $R$ la sous-liste de $L$ entre i et j inclus
		}
	}
	\Retour{R}
\end{algorithm}

\begin{proposition}
	L'algorithme naïf de la décidabilité réaliste (celui naïf de la décidabilité où l'on remplace $Parties(\V)$ par $R\acute{e}alistes(\V, \preccurlyeq)$) a une complexité temporelle de $O(N^3)$.
\end{proposition}

\begin{proof}
	La complexité temporelle de l'extraction d'une sous-liste est linéaire en la taille de la liste. Donc l'algorithme calculant l'ensemble des coalitions réalistes a une complexité en $\Theta(N^3)$.
	
	De plus, $\# \R \leqslant N^2$.
	Dans l'algorithme 3, le corps des deux boucles se font en $O(N)$.
	Ainsi la complexité de l'algorithme naïf de la décidabilité réaliste est en $O(N^3)$.
\end{proof}

Ainsi $B_i^r$ et $B_i'^r$ peuvent être calculés avec un complexité temporelle en $O(N^3)$.